\documentclass[11pt]{article}

\usepackage[margin=1in]{geometry}
\usepackage{amsmath,amssymb,amsthm}
\usepackage{bm}
\usepackage{booktabs}
\usepackage{graphicx}
\usepackage{siunitx}
\usepackage[round,numbers]{natbib}
\usepackage[colorlinks=true,linkcolor=blue,citecolor=blue,urlcolor=blue]{hyperref}

% ---------------------------------------------------------------- shortcuts
\newcommand{\R}{\mathbb{R}}
\newcommand{\pdd}[2]{\frac{\partial #1}{\partial #2}}
\newcommand{\pddt}[2]{\frac{\partial^{2} #1}{\partial #2^{2}}}
\newcommand{\erfc}{\operatorname{erfc}}
\newcommand{\Lpde}{\mathcal{L}_{\mathrm{PDE}}}
\newcommand{\Lic}{\mathcal{L}_{\mathrm{IC}}}
\newcommand{\Lbc}{\mathcal{L}_{\mathrm{BC}}}
\newcommand{\Ldata}{\mathcal{L}_{\mathrm{data}}}
\newcommand{\norm}[1]{\left\lVert #1 \right\rVert}

\title{Phase 1 Sandbox: One-Dimensional Advection--Dispersion Equation\\
\large Mathematical formulation for PINN and FNO verification}
\author{Sergio Lucero}
\date{\today}

\begin{document}
\maketitle

\begin{abstract}
This document fixes the mathematical problem used to debug the physics-informed
neural network (PINN) and Fourier neural operator (FNO) models. Following
Phase~1 of the debugging plan, we strip the governing equation to one spatial
dimension, remove reaction and heterogeneity, and compare the learned solution
against the closed-form Ogata--Banks (1961) solution. Everything below is
non-dimensionalized so that the residual is well scaled (this also anticipates
the variable-scaling check in Phase~2).
\end{abstract}

\section{Governing equation and assumptions}
\label{sec:governing}
The starting point is the advection--dispersion equation (ADE) for a
conservative (non-reactive, non-sorbing) solute in a one-dimensional,
homogeneous, isotropic aquifer with steady, uniform flow:

\begin{equation}
\label{eq:ade_dim}
\pdd{C}{t} + v\,\pdd{C}{x} \;=\; D\,\pddt{C}{x},
\qquad (x,t)\in(0,L)\times(0,T],
\end{equation}
where
\begin{itemize}
  \item $C(x,t)$ \;[$\si{M.L^{-3}}$] is the dissolved solute concentration,
  \item $v$ \;[$\si{L.T^{-1}}$] is the constant average linear (pore) velocity,
  \item $D$ \;[$\si{L^{2}.T^{-1}}$] is the longitudinal hydrodynamic dispersion
        coefficient $D=\alpha_L v + D_m$, with $\alpha_L$ the longitudinal
        dispersivity and $D_m$ the effective molecular diffusion.
\end{itemize}

\paragraph{Phase 1 simplifications.} Compared with the full porous-media model:
\begin{enumerate}
  \item \textbf{One dimension only} --- transverse dispersion is discarded.
  \item \textbf{Homogeneous \& isotropic} --- $v$ and $D$ are constant (no
        heterogeneous conductivity field, no spatially varying retardation).
  \item \textbf{No reaction / no sorption} --- the retardation factor $R=1$ and
        the first-order decay rate $\lambda=0$, hence no source terms.
\end{enumerate}

\section{Initial and boundary conditions}
\label{sec:bc}

We model a clean aquifer invaded from the left by a step (continuous) injection
of concentration $C_0$:

\begin{subequations}
\label{eq:icbc}
\begin{align}
C(x,0) &= 0, && x\in[0,L], \label{eq:ic}\\[2pt]
C(0,t) &= C_0, && t>0, \label{eq:bc_in}\\[2pt]
C(L,t) &= 0, && t>0. \label{eq:bc_out}
\end{align}
\end{subequations}
Equation~\eqref{eq:bc_in} is a Dirichlet (first-type) inlet boundary and
\eqref{eq:bc_out} a Dirichlet ``clean'' outlet. The outlet condition is exact in
the limit of a semi-infinite domain; in the sandbox we choose $L$ (equivalently
the final time $T$) large enough that the plume front stays well inside the
domain, so that $C(L,t)\approx 0$ and the semi-infinite reference solution of
Section~\ref{sec:analytical} applies to machine precision. The residual
approximation error is quantified in Section~\ref{sec:verification}.

\section{Non-dimensionalization}
\label{sec:nondim1d}

Introduce the dimensionless variables
\begin{equation}
\xi=\frac{x}{L},\qquad
\tau=\frac{t\,v}{L}\quad(\text{pore volumes}),\qquad
u=\frac{C}{C_0},
\end{equation}
and the \textbf{P\'eclet number}
\begin{equation}
\label{eq:pe}
\mathrm{Pe}=\frac{vL}{D},
\end{equation}
which measures advective versus dispersive transport. Substitution into
\eqref{eq:ade_dim} gives the dimensionless ADE
\begin{equation}
\label{eq:ade_nd}
\boxed{\;
\pdd{u}{\tau}+\pdd{u}{\xi}=\frac{1}{\mathrm{Pe}}\pddt{u}{\xi},
\;}
\qquad (\xi,\tau)\in(0,1)\times(0,\tau_{\max}],
\end{equation}
with $\tau_{\max}=Tv/L$, subject to
\begin{equation}
u(\xi,0)=0,\qquad u(0,\tau)=1,\qquad u(1,\tau)=0 .
\label{eq:nd_bc}
\end{equation}
All quantities in \eqref{eq:ade_nd}--\eqref{eq:nd_bc} are $\mathcal{O}(1)$,
which keeps the PINN residual well conditioned.

\section{Analytical (reference) solution}
\label{sec:analytical}

For the semi-infinite problem with a step inlet, the classical solution
\citep{ogata1961solution} is
\begin{equation}
\label{eq:ogata_nd}
u_{\mathrm{ref}}(\xi,\tau)=
\frac{1}{2}\left[
\erfc\!\left(\frac{(\xi-\tau)\sqrt{\mathrm{Pe}}}{2\sqrt{\tau}}\right)
+\;e^{\mathrm{Pe}\,\xi}\,
\erfc\!\left(\frac{(\xi+\tau)\sqrt{\mathrm{Pe}}}{2\sqrt{\tau}}\right)
\right],
\end{equation}
or, in dimensional form,
\begin{equation}
\label{eq:ogata_dim}
C_{\mathrm{ref}}(x,t)=
\frac{C_0}{2}\left[
\erfc\!\left(\frac{x-vt}{2\sqrt{Dt}}\right)
+\;e^{vx/D}\,
\erfc\!\left(\frac{x+vt}{2\sqrt{Dt}}\right)
\right].
\end{equation}
Equation~\eqref{eq:ogata_nd} is the ground truth used for the Phase~1 baseline
test. Note that $u_{\mathrm{ref}}(0,\tau)=1$ and
$\lim_{\xi\to\infty}u_{\mathrm{ref}}=0$, so it satisfies
\eqref{eq:bc_in} exactly and \eqref{eq:bc_out} in the far-field limit.

\section{PINN formulation}
\label{sec:pinn}

We approximate $u\approx u_\theta(\xi,\tau)$ with a feed-forward multilayer
perceptron (MLP) $\theta\in\R^{n_\theta}$, with $\tanh$ activations in the hidden
layers and a logistic (sigmoid) output so that the prediction is hard-bounded to
$u\in(0,1)$, consistent with the normalized concentration. The PDE residual is
\begin{equation}
\label{eq:residual}
\mathcal{R}_\theta(\xi,\tau)
=\pdd{u_\theta}{\tau}+\pdd{u_\theta}{\xi}
-\frac{1}{\mathrm{Pe}}\pddt{u_\theta}{\xi},
\end{equation}
computed exactly by automatic differentiation. The training objective is the
composite loss
\begin{equation}
\label{eq:loss}
\mathcal{L}(\theta)=
\lambda_{\mathrm{PDE}}\,\Lpde(\theta)
+\lambda_{\mathrm{IC}}\,\Lic(\theta)
+\lambda_{\mathrm{BC}}\,\Lbc(\theta)
+\lambda_{\mathrm{data}}\,\Ldata(\theta),
\end{equation}
with
\begin{align}
\Lpde(\theta)&=\frac{1}{N_f}\sum_{i=1}^{N_f}
\bigl|\mathcal{R}_\theta(\xi_i,\tau_i)\bigr|^{2},
&&(\xi_i,\tau_i)\sim\mathcal{U}\bigl((0,1)\times(0,\tau_{\max})\bigr),
\label{eq:Lpde}\\[4pt]
\Lic(\theta)&=\frac{1}{N_{\mathrm{ic}}}\sum_{i=1}^{N_{\mathrm{ic}}}
\bigl|u_\theta(\xi_i,0)-0\bigr|^{2},
&&\xi_i\sim\mathcal{U}(0,1),
\label{eq:Lic}\\[4pt]
\Lbc(\theta)&=\frac{1}{N_{\mathrm{bc}}}\sum_{i=1}^{N_{\mathrm{bc}}}
\Bigl(\bigl|u_\theta(0,\tau_i)-1\bigr|^{2}
+\bigl|u_\theta(1,\tau_i)-0\bigr|^{2}\Bigr),
&&\tau_i\sim\mathcal{U}(0,\tau_{\max}).
\label{eq:Lbc}
\end{align}
\paragraph{Soft versus hard boundary conditions (Phase 2).}
In \eqref{eq:Lbc} the Dirichlet data enter as a penalty, which is a \emph{soft}
constraint: the boundary error is only driven to zero as training proceeds, and
the PDE and boundary terms compete for gradient budget. Phase~2 therefore moves
to a \emph{hard} (architectural) ansatz that satisfies the initial condition and
the outlet condition exactly and the inlet condition up to a short ramp,
\begin{equation}
\label{eq:hardbc}
u_\theta(\xi,\tau)=(1-\xi)\,s(\tau)
+\xi(1-\xi)\,s(\tau)\,\mathrm{NN}_\theta(\xi,\tau),
\qquad s(\tau)=1-e^{-\tau/\varepsilon},
\end{equation}
with $\varepsilon=10^{-3}$ (Phase~2 found $\varepsilon=10^{-4}$ unstable: too
steep a ramp near $\tau=0$). Both terms of \eqref{eq:hardbc} carry the factor
$s(\tau)$, so $u_\theta(\xi,0)=0$ for every $\xi$: the initial condition
\eqref{eq:ic} holds identically. At $\xi=1$ the factor $(1-\xi)$ vanishes in both
terms, so the outlet condition \eqref{eq:bc_out} also holds identically. Only the
inlet is approximate: $u_\theta(0,\tau)=s(\tau)$, which equals $1$ only in the
limit $\tau\gg\varepsilon$, recovering \eqref{eq:bc_in} away from the corner.

\paragraph{A structural floor on the residual norm.}
The ramp in \eqref{eq:hardbc} has a consequence that matters when reading
residual diagnostics: since $s'(\tau)=e^{-\tau/\varepsilon}/\varepsilon$, we
have $s'(0)=1/\varepsilon$, so
\begin{equation}
\label{eq:residual_floor}
\max_{\xi,\tau}\bigl|\mathcal{R}_\theta\bigr|
\;\gtrsim\;\frac{1}{\varepsilon}
\end{equation}
\emph{independently of the network weights}. Measured at the entrance corner,
$|R|$ equals $1/\varepsilon$ to three digits at initialization and after
training, and for every seed. This is a property of the ansatz, not a training
failure, and no amount of optimization removes it. It is the residual-side
analogue of the corner discontinuity in Section~\ref{sec:verification}: both say
that a small set of points is structurally excluded from the error metric, which
is why band-limited means (interior, inlet, inlet-adjacent) are reported
instead of a global maximum. The optional data term $\Ldata$ supervises against
\eqref{eq:ogata_nd} at a few interior points when used as a sanity check, and is
set to $\lambda_{\mathrm{data}}=0$ for the pure-physics baseline.

\paragraph{Optimization.} Adam is used for the first $\sim\!10^4$--$2\times10^4$
iterations (with collocation points re-sampled each step), followed by a short
L-BFGS polish; this ordering is motivated in Phase~3 of the plan.

\section{FNO formulation}
\label{sec:fno}

The FNO learns the \emph{solution operator} rather than a single trajectory.
Writing the governing parameters as an input function $a$ and the concentration
field as the output $u$, the operator to learn is
\begin{equation}
\mathcal{G}: a \;\longmapsto\; u,
\qquad a=(\mathrm{Pe})\ \text{or}\ (v,D),
\qquad u=u(\xi,\tau),
\end{equation}
discretized as
\begin{equation}
u_{j}^{n}\approx
\mathcal{G}_\theta\!\bigl(a;\ \xi_j,\tau_n\bigr),\qquad
j=0,\dots,N_x,\quad n=0,\dots,N_t .
\end{equation}
$\mathcal{G}_\theta$ is a Fourier neural operator
\citep{li2020fourier}: a lifting MLP, $K$ spectral (Fourier) convolution layers
in which the integral kernel is parameterized directly in frequency space with
$k_{\max}$ retained modes, and a projection MLP. Phase~1 uses the synthetic
dataset of Section~\ref{sec:data} (a family of ``hydrologies'' indexed by
$(v,D)$) to train the operator in supervised mode against
\eqref{eq:ogata_nd}; spectral-hyperparameter checks (number of modes, resolution
invariance) are deferred to Phase~4.

\section{Synthetic data generation}
\label{sec:data}

The reference generator produces a family of 1D transport scenarios
(``hydrologies'') by sampling the physical parameters and evaluating
\eqref{eq:ogata_dim} on a space--time grid:

\begin{itemize}
  \item \textbf{Scenario:} a tuple $(C_0,v,D,L,T)$ with fixed grid
        $(N_x,N_t)$.
  \item \textbf{Single-scenario output:} the full field
        $C(x_i,t_n)$ plus the final-time snapshot $C(x_i,T)$.
  \item \textbf{Dataset:} $N$ scenarios with $(v,D)$ drawn from prescribed
        ranges; each sample stores the parameters and the corresponding
        concentration field, which is exactly the supervised input/output format
        required by the FNO in Section~\ref{sec:fno}.
\end{itemize}

\section{Verification metric}
\label{sec:verification}

The baseline passes if the PINN converges to the reference solution. The
reported metric is the relative $L^2$ error over the domain,
\begin{equation}
\label{eq:metric}
E_{\mathrm{rel}}
=\frac{\norm{u_\theta-u_{\mathrm{ref}}}_{2}}{\norm{u_{\mathrm{ref}}}_{2}},
\qquad
\norm{f}_{2}^{2}=\frac{1}{N_x N_t}
\sum_{j,n} f(\xi_j,\tau_n)^{2},
\end{equation}
evaluated on a fixed test grid independent of the training collocation points.
A target of $E_{\mathrm{rel}}\lesssim 10^{-2}$ is the scale at which the core
architecture and loss formulation are confirmed; a much larger error indicates
a bug in the residual/graph rather than in heterogeneity or dimensionality
(Phase~2 onward).

\paragraph{Reporting policy: medians, not best runs.}
Because training is stochastic, $E_{\mathrm{rel}}$ is a random variable. We
therefore report its \emph{median} over $N\geq5$ seeds together with the
interquartile range, and never a single run in isolation. Section~\ref{sec:results}
shows why this is not pedantry: at fixed hyperparameters the metric spans more
than an order of magnitude across seeds, so a single-run comparison can invert
the sign of a real effect. A candidate variant counts as an improvement only if
it beats the baseline median by more than the baseline IQR, evaluated on the
same seeds (a paired comparison).

\paragraph{Corner singularity and the interior metric.}
The step-inlet problem is discontinuous at the entrance corner
$(\xi,\tau)=(0,0)$: the initial condition \eqref{eq:ic} imposes $u=0$ while the
inlet boundary \eqref{eq:bc_in} imposes $u=1$. The reference solution
\eqref{eq:ogata_nd} returns $u_{\mathrm{ref}}(0,0)=1$, whereas any continuous
ansatz that honors the initial condition returns $0$ there. The pointwise error
at the corner is therefore exactly $1$ and a thin band around it is
contaminated; this is a metric artifact of the discontinuous problem, not a
network or loss bug. We therefore always report a second, corner-excluded
metric alongside the full-grid one,
\begin{equation}
\label{eq:metric_int}
E_{\mathrm{rel}}^{\mathrm{int}}
=\frac{\norm{u_\theta-u_{\mathrm{ref}}}_{2,\Omega_{\mathrm{int}}}}
       {\norm{u_{\mathrm{ref}}}_{2,\Omega_{\mathrm{int}}}},\qquad
\Omega_{\mathrm{int}}=\{(\xi,\tau): \xi>0.05,\ \tau>0.05\},
\end{equation}
which is the physically meaningful number for the interior plume.

\section{Baseline results}
\label{sec:results}

The Phase~1 hard-BC ansatz (Section~\ref{sec:pinn})
$u_\theta(\xi,\tau)=(1-\xi)s(\tau)+\xi(1-\xi)s(\tau)\,\mathrm{NN}_\theta(\xi,\tau)$
with $s(\tau)=1-e^{-\tau/\varepsilon}$, $\varepsilon=10^{-3}$.

\paragraph{Seed variance dominates a single run.}
Training is stochastic (collocation resampling and, most importantly, the
initialization), so a single run does not characterize the method. Repeating
the baseline over five seeds ($0$--$4$, identical hyperparameters) gives

\begin{table}[h]
\centering
\begin{tabular}{lccccc}
\toprule
 & median & IQR & min & max & spread \\
\midrule
$E_{\mathrm{rel}}$ (full grid)      & $1.61\times10^{-2}$ & $2.18\times10^{-2}$ & $1.17\times10^{-2}$ & $8.14\times10^{-2}$ & $7.0\times$ \\
$E_{\mathrm{rel}}^{\mathrm{int}}$ ($\Omega_{\mathrm{int}}$) & $1.42\times10^{-2}$ & $2.59\times10^{-2}$ & $6.64\times10^{-3}$ & $8.16\times10^{-2}$ & $12.3\times$ \\
\bottomrule
\end{tabular}
\caption{Phase~1 hard-BC baseline over five seeds. ``Spread'' is
$\max/\min$. The distribution is right-skewed: the mean of the full-grid metric
($3.19\times10^{-2}$) is twice its median, so median and IQR are reported
rather than mean and standard deviation. Reproduce with
\texttt{runs/sweeps/baseline\_n5} (harness: \texttt{SCRIPTS/seed\_sweep.py}).}
\label{tab:baseline_seeds}
\end{table}

Two conclusions follow, and they supersede the single-run figures quoted in
earlier drafts of this document.

First, \textbf{the corner-excluded metric does not rescue the target}. The best
seed reaches $E_{\mathrm{rel}}^{\mathrm{int}}=6.64\times10^{-3}$, comfortably
below $10^{-2}$; the median ($1.42\times10^{-2}$) is not. Reporting the best run
alone would therefore overstate the method. The measure-zero corner
discontinuity of Section~\ref{sec:verification} is real and accounts for part of
the gap between the two metrics, but it is not the whole story. Across the
five seeds the two metrics are perfectly rank-correlated (Spearman
$\rho=1$), so excluding the corner never reorders the runs, and the ratio
$E_{\mathrm{rel}}^{\mathrm{int}}/E_{\mathrm{rel}}$ ranges from $0.57$ to
$1.00$ --- on the worst seed the interior restriction removes essentially no
error. Even in the interior the error therefore varies by more than an order of
magnitude ($12.3\times$); the corner is a contributing artifact, not the
governing one.

Second, \textbf{run-to-run variance is now the dominant error source}, ahead of
the loss-shaping refinements of Phase~3. A variant must beat the median by more
than the IQR of
Table~\ref{tab:baseline_seeds} to be considered an improvement. As a concrete
example, causal $\tau$-ramped PDE weighting looked like a $32\%$ gain on seed
$0$ but won only $2$ of $5$ paired seeds (median ratio $1.31$ in the wrong
direction), so it is not adopted. All subsequent claims in this project are
reported as a median over $N\geq5$ seeds with the spread alongside.

\section{FNO results: five seeds, and the modes that matter}
\label{sec:fno-results}

Section~\ref{sec:fno} deferred the spectral-hyperparameter checks --- number of
retained modes, resolution transfer, epoch budget --- to Phase~4. They have now
been run, over five seeds each, and the first thing they produced was a defect
worth recording.

\paragraph{The operator that could only return a constant.}
The first FNO run reported $E_{\mathrm{rel}}=0.973$. Because
\eqref{eq:metric} is normalized by $\norm{u_{\mathrm{ref}}}_{2}$, the predictor
$u\equiv 0$ scores exactly $1$: the network was \emph{worse than predicting
nothing}, and only a factor of three better than the mean-field reference
($0.221$) that a one-parameter fit achieves. The cause was structural, not
optimizational. The input was two spatially constant channels $(\mathrm{Pe})$
or $(v,D)$; the lifting layer is pointwise, and a spectral convolution
multiplies Fourier coefficients, so a constant input populates only the DC mode
and every layer maps a constant to a constant. The measured spatial standard
deviation of the output was $2\times10^{-10}$, and a profile-only error of
$1.000$ confirmed it: the architecture had no way to represent $\xi$ or $\tau$
at all. Appending normalized $(\xi,\tau)$ coordinate channels as inputs fixed
it, and \texttt{solve\_1d\_fno.py} now warns when a run fails to beat both the
all-zeros and the mean-field predictor, so the failure cannot pass silently
again. The general lesson is the one \eqref{eq:metric} invites: a surrogate
scored on a relative $L^{2}$ error must be reported next to the trivial
predictors, because a normalized error can look small while the operator is
returning nothing.

\paragraph{Seed $0$ is not the method.}
With the coordinate channels in place the operator learns. Table
\ref{tab:fno-seeds} reports the 250-epoch and 120-epoch configurations over
seeds $0$--$4$ with identical hyperparameters otherwise.

\begin{table}[h]
\centering
\begin{tabular}{lcccc}
\toprule
configuration & median & IQR & min & max \\
\midrule
$250$ epochs & $1.99\times10^{-2}$ & $4.60\times10^{-3}$ & $1.46\times10^{-2}$ & $4.59\times10^{-2}$ \\
$120$ epochs & $3.02\times10^{-2}$ & $3.17\times10^{-3}$ & $2.35\times10^{-2}$ & $5.98\times10^{-2}$ \\
mean field, $1$ parameter & $2.21\times10^{-1}$ & --- & --- & --- \\
zero predictor & $1.00$ & --- & --- & --- \\
\bottomrule
\end{tabular}
\caption{FNO test error over five seeds (coordinate channels on, $12$ retained
modes, $128\times128\times32$ scenarios), with the two trivial predictors for
scale; $\max/\min$ spread is $3.1\times$ and $2.6\times$ for the two budgets.
The figure first quoted for this operator, $1.46\times10^{-2}$, is seed $0$: the
most favourable of the five draws. The honest statement is the median,
$1.99\times10^{-2}$ at $250$ epochs. Reproduce with
\texttt{runs/sweeps/fno\_coords\_250ep} and
\texttt{runs/sweeps/fno\_coords\_120ep}.}
\label{tab:fno-seeds}
\end{table}

Compared seed by seed, $250$ epochs wins on all five (error ratios
$0.62$, $0.66$, $0.60$, $0.68$, $0.77$; median $0.66$), so the extra budget is
a real improvement and not a lucky ordering. Against the mean-field reference
the median is $11.1$ times better; at seed $0$ it is $15.1$ times better, which
is where the ``$15\times$'' figure that circulated came from, and it should be
quoted as $11\times$ for the same reason the median is quoted at all.

\paragraph{Retained modes.}
Table~\ref{tab:fno-modes} sweeps $k_{\max}\in\{4,12,24\}$ at the 120-epoch
budget, so that the comparison is affordable over five seeds.

\begin{table}[h]
\centering
\begin{tabular}{lcccc}
\toprule
$k_{\max}$ & median & IQR & min & max \\
\midrule
$4$  & $4.06\times10^{-2}$ & $4.48\times10^{-3}$ & $2.79\times10^{-2}$ & $7.18\times10^{-2}$ \\
$12$ & $3.02\times10^{-2}$ & $3.17\times10^{-3}$ & $2.35\times10^{-2}$ & $5.98\times10^{-2}$ \\
$24$ & $2.47\times10^{-2}$ & $1.72\times10^{-3}$ & $2.19\times10^{-2}$ & $5.00\times10^{-2}$ \\
\bottomrule
\end{tabular}
\caption{Truncation sweep at $120$ epochs over five seeds. Medians decrease
monotonically with $k_{\max}$ and the spread narrows; paired by seed, $k_{\max}=24$
beats $k_{\max}=4$ on all five (median ratio $0.66$). Reproduce with
\texttt{runs/sweeps/fno\_modes\_120ep}.}
\label{tab:fno-modes}
\end{table}

The default $k_{\max}=12$ is therefore not at the knee: the operator is still
improving at $24$ modes, and the residual error there
($2.47\times10^{-2}$) is dominated by the epoch budget, not by truncation. As
with the epoch comparison, every single-seed value in this sweep was optimistic
relative to its median, which is the argument for running these sweeps at all.

\paragraph{What is still a single draw.}
Three rows remain preliminary by the standard above, and are quoted as such:
the resolution-transfer checks ($64^{2}$: $4.14\times10^{-2}$,
$32^{2}$: $6.92\times10^{-2}$) and the profile-only error
($7.1\times10^{-2}$) were each measured on one seed, and the
$300$-epoch plateau has one draw. None of them supports a claim stronger than
``consistent with'', and none is used in Chapter~\ref{chap:conclusions}.

\section{Baseline parameters}
\label{sec:params}

\begin{table}[h]
\centering
\begin{tabular}{llll}
\toprule
Symbol & Meaning & Value & Units \\
\midrule
$L$   & domain length        & $100$            & \si{m} \\
$v$   & pore velocity        & $0.1$            & \si{m.d^{-1}} \\
$D$   & dispersion           & $0.5$            & \si{m^{2}.d^{-1}} \\
$C_0$ & inlet concentration  & $1$              & \si{mg.L^{-1}} \\
$\mathrm{Pe}=vL/D$ & P\'eclet number & $20$ & -- \\
$\tau_{\max}=Tv/L$ & final pore volumes & $0.5$ & -- \\
\midrule
$N_x,N_t$ & reference grid & $256\times 256$ & -- \\
$N_f$   & PDE collocation points (per step) & $1500$ & -- \\
$N_{\mathrm{ic}}$ & IC points & $150$ & -- \\
$N_{\mathrm{bc}}$ & BC points & $150$ & -- \\
\bottomrule
\end{tabular}
\caption{Phase 1 sandbox configuration. The outlet BC
\eqref{eq:bc_out} is consistent with the semi-infinite reference solution
because $\tau_{\max}=0.5$ keeps the front (centered at $\xi=\tau$) well inside
the domain.}
\label{tab:params}
\end{table}

\bibliographystyle{plainnat}
\begin{thebibliography}{9}
\bibitem{ogata1961solution}
A. Ogata and R. B. Banks,
\textit{A solution of the differential equation of longitudinal dispersion in
porous media}, U.S. Geological Survey Professional Paper 411-A, 1961.
\bibitem{li2020fourier}
Z. Li, N. Kovachki, K. Azizzadenesheli, B. Liu, K. Bhattacharya, A. Stuart,
A. Anandkumar, \textit{Fourier Neural Operator for Parametric Partial
Differential Equations}, ICLR, 2021.
\end{thebibliography}

\end{document}
